\documentclass[11pt,letterpaper]{article}
\usepackage[T1]{fontenc}
\usepackage{lmodern,microtype,amsmath,amssymb,amsthm,booktabs}
\usepackage[margin=0.95in,headheight=14pt]{geometry}
\usepackage{fancyhdr,xcolor}
\usepackage[colorlinks=true,linkcolor=black,urlcolor=blue!55!black]{hyperref}
\newtheorem{theorem}{Theorem}
\newtheorem{lemma}[theorem]{Lemma}
\newtheorem{corollary}[theorem]{Corollary}
\newcommand{\oeis}[1]{\href{https://oeis.org/A#1}{A#1}}
\setlength{\parskip}{0.18em}
\raggedbottom
\hypersetup{pdftitle={Plateaus in the polygonal partition array},pdfauthor={Blake Stone, with OpenAI Codex}}
\pagestyle{fancy}\fancyhf{}\lhead{\small Plateaus in the polygonal partition array}\rhead{\small Blake Stone}\cfoot{\thepage}
\begin{document}
\thispagestyle{plain}
\begin{center}
{\LARGE\bfseries All plateaus in the polygonal partition array\par}
\vspace{0.5em}
Blake Stone\\
{\small Research manuscript prepared with OpenAI Codex}
\end{center}

\begin{abstract}
Omar E. Pol conjectured the lengths and levels of every constant run in the columns of OEIS \oeis{195825}. We determine their positions as well. In column $k$ the plateau indexed by $r$ starts at $r(k+3)$, has length $k+1-2r$, and has level equal to the number of three-colored partitions of integers at most $r$. These are all the plateaus of length at least two.
\end{abstract}

\section{From polygonal numbers to partitions}
Fix $k\ge1$, set $m=k+2$, and write $C_m(n)$ for the entry in row $n\ge0$ of the column. Its defining recurrence in \oeis{195825} uses the nonzero generalized polygonal numbers $P(j)=mj(j-1)/2+j$ for $j=1,-1,2,-2,\ldots$, with recurrence signs $+,+,-,-,\ldots$. Equivalently its generating function is the reciprocal of $\sum_{j\in\mathbb Z}(-1)^j q^{P(j)}$. Jacobi's triple product with base $q^m$ and argument $q$ gives
\[
 \sum_{n\ge0}C_m(n)q^n
  =\prod_{\substack{b\ge1\\ b\equiv0,1,-1\;(\mathrm{mod}\ m)}}(1-q^b)^{-1}.
\]
The three residues are distinct even at $m=3$. Thus the column counts partitions into these parts. This also identifies the proposed columns $k=2,3$ with \oeis{006950} and \oeis{036820}. For the first identification, replace $1+q^a$ by $(1-q^{2a})/(1-q^a)$ for every odd $a$ in the product for \oeis{006950}. For the second, the two forbidden residues modulo five are precisely $2$ and $3$.

\section{Where the counts increase}
Let $b_m(n)=C_m(n)-C_m(n-1)$ with $C_m(-1)=0$. Removing one part $1$ from partitions containing one shows that $b_m(n)$ counts partitions of $n$ with no part $1$. Their possible parts are $tm+e$, $t\ge1$, $e=-1,0,1$. Every such part lies in the additive semigroup generated by $m-1,m,m+1$, since $tm+e=(t-1)m+(m+e)$. Conversely those three generators are allowed. Sums of exactly $r$ generators fill the entire integer interval
\[
                       I_r=[r(m-1),r(m+1)] .
\]
Indeed the change from $rm$ can be any integer between $-r$ and $r$. Including $I_0=\{0\}$, the union of these intervals is exactly the support of $b_m$.

\begin{theorem}
The maximal constant runs of length at least two have indices
$0\le r\le\lfloor(k-1)/2\rfloor$. Their first and last row indices are
\[
                  r(m+1),\qquad (r+1)(m-1)-1.
\]
Their lengths are $k+1-2r$. At the $r$th such run the level is
\[
                 L_r=\sum_{s=0}^r p_3(s),\qquad
      \sum_{s\ge0}p_3(s)z^s=\prod_{t\ge1}(1-z^t)^{-3}.
\]
\end{theorem}
\begin{proof}
Between $I_r$ and $I_{r+1}$ there is a missing integer exactly when $m-2r-2\ge1$. Both bounding endpoints are in the support. Hence the column stays constant from the last point of $I_r$ to the integer preceding $I_{r+1}$, giving the stated position and length. The intervals have increasing endpoints and their gaps shrink by two, so their union has no later gaps.

For the level, regard the part $tm+e$ as a part of weight $t$ and color $e$. Its representation is unique because $m\ge3$. A colored partition of total weight $s$ has actual size in $I_s$, as its sum of colors is between $-s$ and $s$. For an $r$ followed by a gap, all $I_s$ with $s\le r$ end at or below $r(m+1)$, and all those with $s>r$ start above it. Therefore summing $b_m(n)$ through $n=r(m+1)$ counts exactly the three-colored partitions of weights at most $r$. This sum equals $C_m(r(m+1))$.
\end{proof}

For instance, $k=6$ gives the intervals of row indices $[0,6]$, $[9,13]$ and $[18,20]$, with levels $1,4,13$. For general $k$ the plateau count is $\lceil k/2\rceil$. At fixed $r$ its level is already attained for every $k\ge2r+1$. The generating function of the levels (starting at $r=0$) is
\[
     \sum_{r\ge0}L_rz^r
       =(1-z)^{-1}\prod_{t\ge1}(1-z^t)^{-3}.
\]
This proves the Euler-transform description asked about in \oeis{210843}: the exponents are $4,3,3,\ldots$. V. Kotesovec already used this product to compute that sequence in 2015.

\begin{thebibliography}{9}
\bibitem{pol} O. E. Pol, OEIS entries \oeis{195825}, \oeis{195152} and \oeis{210843}; plateau conjecture, 2012. See also the individual columns \oeis{006950} and \oeis{036820}.
\bibitem{dlmf} NIST Digital Library of Mathematical Functions, \href{https://dlmf.nist.gov/17.8.E1}{equation 17.8.1}, Jacobi's triple product.
\end{thebibliography}
\end{document}
